% ===========================================================================
%  01 — 引言
%
%  事实来源：paper_context/paper_outline.md §2.1
%  证据：     paper_context/proposal_summary.md
%             paper_context/system_delta_from_proposal.md
%             paper_context/claims_registry.md
%  术语：     paper_context/terminology.md
%  声明：     C-PROPOSAL-01, C-PROPOSAL-02（明确标注为愿景，而非已完成）
% ===========================================================================
\chapter{引言}\label{ch:introduction}

本文描述了 SimTutor 的设计、实现与技术评估。SimTutor 是一个基于遥测数据基础的教学运行时系统，结合了程序包驱动的状态跟踪、检索增强的语言模型辅助生成以及微调视觉语言模型（VLM）的视觉事实提取功能，以支持高保真飞行仿真器中的程序性学习。本研究以 DCS World 中 F/A-18C 冷启动任务为代表案例，针对复杂、序列敏感的程序性教学挑战展开，同时弥合了原始研究愿景（基于虚拟现实的教学系统）与截至目前已完成的技术基线之间的差距。本章定义研究问题，追溯从研究计划到当前实现的演进过程，列举本研究的贡献，并概述论文剩余部分的组织结构。

% ===========================================================================
\section{研究动机}\label{sec:intro-motivation}

程序性学习——获取必须在正确上下文中正确执行的固定有序动作序列——面临着与其区别于陈述性学习或概念性学习的独特挑战。学习者不仅需要记忆步骤，还必须识别每个步骤的适用时机，验证其前提条件是否满足，在推进前确认其完成，并在出错时恢复，而不将错误传播到程序的后续部分。在高保真仿真环境中，这些挑战因座舱或工作空间的复杂性、仪表信息的密度以及在执行不熟悉序列时保持态势感知的认知负荷而进一步加剧。

传统的培训辅助手段——视频演示、纸质检查单、教员引导示范——在初始熟悉阶段有效，但会打断仿真器内练习的连续性。暂停练习查阅手册或重播视频片段的学习者必须从座舱环境中切换注意力，从而丧失沉浸感，在某些情况下甚至会丢失最初触发查阅需求的瞬态仿真器状态。相反，来自大语言模型（LLM）的缺乏系统状态基础的纯文本辅助可能提供流畅但事实性错误的指导，因为模型无法获取当前仿真器状态。在这两个极端之间，需要一种\emph{基于状态感知证据的、上下文相关}的教学方式：一个通过实时仿真器遥测数据观察学习者操作、检索相关文档、在遥测数据不足时从座舱中提取视觉证据、并在不中断训练流程的前提下提供简洁、可验证指导的系统。

DCS World 中 F/A-18C 的冷启动程序是本论文的训练场景。该程序包含 25 个步骤（S01--S25），分为六个阶段（P1--P6），从电瓶供电开始，历经发动机启动、航电配置、飞行控制检查，直至滑行准备（详见\S\ref{sec:bg-coldstart}）。每个步骤都带有显式的前提条件，这些条件取决于先前步骤的完成情况以及特定遥测条件的满足程度：例如，左发起动（S10）要求右发参数——转速、温度、燃油流量、尾喷口位置和滑油压力——在启动左发之前处于标称范围内。这些步骤间依赖关系使得冷启动任务成为航空、海事操作、工业维护和医疗程序指导中基于仿真器培训的序列敏感、状态依赖型程序性任务这一更广泛类别的代表性实例。TODO:CITE —— 智能教学系统和程序性学习文献。

本研究的核心前提是：通过将每个可操作的输出都建立于显式证据之上——遥测变量、门控规则评估、检索到的文档片段或视觉事实观察——可以使教学系统显著提升可靠性与可审计性。这种证据溯源约束不仅仅是一种架构偏好，更是一项安全要求：在座舱程序的语境中，被错误高亮的控件或被错误识别的程序状态可能会使学习者产生混淆或误导。本文所描述的 SimTutor 系统在一个包含自动化测试、基准测试工件和回放基础设施的可运行运行时环境中实例化了这一原则。

% ===========================================================================
\section{论文愿景与范围}\label{sec:intro-vision}

本论文的原始研究愿景，如研究计划所述，核心是在 Meta Quest~3 上结合 DCS-VR~\citep{ThesisProposal1128} 构建一个以虚拟现实为先导的基于状态证据的教学系统。计划的系统将通过 VR 座舱内的眼动追踪、手部追踪和语音输入接收学习者交互；检索相关手册和检查单摘录；通过检索增强的 LLM 生成简洁、附带引文依据的步骤指导；并将指导渲染为头显内步骤卡片或注册到座舱控件上的叠加标注。研究计划进一步规定了一个三条件受控人因评估，比较视频加笔记基线、无状态基础的纯文本 LLM 以及基于状态证据的教学系统，评估指标包括完成率、程序性错误率、任务时间、NASA-TLX 工作负荷评分、前/后知识测验，以及在 2--4 周间隔后进行的保持和迁移测试。

该愿景定义了本论文的长期研究议程。然而，截至目前已完成的工作代表了技术基线，而非完整 VR 教学系统的实现。当前实现在桌面/DCS 运行时上运行，使用座舱叠加高亮而非头显内步骤卡片作为主要输出模态，并通过自动化基准测试而非人因实验进行验证。VR 部署和人因评估仍为计划的下一阶段（见第\ref{ch:evaluation}章，第\ref{sec:eval-planned}节）。

我们在开篇即明确记录这一差距，以在原始研究愿景与已完成的技术贡献之间建立清晰的边界。本论文应被理解为：一项系统与方法论贡献，奠定了技术基础——运行时、数据采集、VLM 适配流水线、基准测试基础设施——VR 和人因评估可在该基础上进行构建。

% ===========================================================================
\section{当前实现范围}\label{sec:intro-scope}

已构建的系统在研究计划的多个重要方面有所不同，而这些差异本身正是本论文贡献的来源之一。详细比较见第\ref{sec:design-delta}节；在此我们总结关键差异以帮助读者定位。

\paragraph{平台。}
当前系统运行于桌面/DCS（非 VR）配置。显示器为标准显示器；输入通过键盘和 HOTAS（手不离杆操纵系统）；输出渲染为 DCS 屏幕上的座舱叠加高亮，而非头显内步骤卡片。Meta Quest~3 VR 适配器尚未实现。

\paragraph{遥测数据基础。}
当前实现包含一个比原计划预期更为精细的遥测子系统：原始 DCS-BIOS UDP 接收、带源丢失跟踪的变量解析、增量去噪处理（阈值、防抖动和逐变量过滤），以及用于提示安全压缩的增量聚合。这些新增功能是由处理嘈杂高频仿真器遥测数据的实际工程需求所驱动的。

\paragraph{VLM 视觉事实提取。}
相对于研究计划最显著的新增内容是 VLM 视觉事实提取流水线。原计划设想一个纯文本 RAG+LLM 导师，以只读仿真器状态作为唯一的系统状态基础信号。在实现过程中，我们发现若干关键程序步骤（S08、S15、S18）需要对未通过 DCS-BIOS 导出的显示页面内容进行视觉确认。VLM 流水线即为弥补这一缺口而添加，并发展成为一个自成一体的研究贡献，涵盖结构化视觉事实本体设计、带人工审查的数据集构建、双语 LoRA 微调，以及跨两个模型骨干的自动化基准测试评估。

\paragraph{程序包和基础设施。}
当前系统将程序定义、门控规则、遥测映射、错误分类法和视觉事实本体组织为声明式 YAML 配置（\emph{程序包}），而非将其隐式编码在提示或代码中。系统还包含事件日志记录与回放基础设施、自动化评分引擎，以及在没有模型可用或模型输出验证失败时运行的确定性降级路径。这些要素在原计划中均未规定。

这些差异并不使原计划失效；它们反映了一个工程现实：基于状态证据的教学运行时系统在可部署于 VR 人因研究之前，需要大量支撑性基础设施。当前系统恰好提供了这一基础设施。

% ===========================================================================
\section{贡献}\label{sec:intro-contributions}

本论文做出以下贡献：

\begin{enumerate}

\item \textbf{程序包驱动、基于遥测数据基础的教学运行时架构。}
我们提出了一个以端口-适配器（六边形）架构组织的教学运行时的设计与实现，将仿真器无关的核心逻辑（程序跟踪、门控、评分、知识检索）与仿真器相关适配器（DCS-BIOS 遥测、叠加渲染、视觉捕获）相分离。系统支持实时帮助循环：运行时从规范化遥测变量、增量摘要、候选程序步骤、确定性步骤提示、叠加目标允许列表、检索到的知识片段以及可选的视觉事实观察中组装结构化上下文；将该上下文传递给 LLM 以生成结构化帮助响应；依据模式、允许列表和证据溯源约束验证响应；执行经过验证的叠加操作，同时记录每个阶段以便下游回放和分析。确定性降级路径确保在没有模型可用或模型输出验证失败时系统仍能正常运行。此贡献在第\ref{ch:system-design}、\ref{ch:methodology}和\ref{ch:implementation}章中描述。
% 证据：C-ARCH-01, C-ARCH-02, C-ARCH-03, C-PROC-01, C-PROC-02, C-TELEM-01, C-RAG-01, C-RUNTIME-01

\item \textbf{结构化视觉事实本体及完整的 VLM 数据、适配与基准测试流水线。}
我们定义了一个包含 13 个视觉事实的本体，将座舱显示页面状态分解为跨三个显示区域（Left DDI、Right DDI、AMPCD）的离散、可审计命题。该本体经历了两代演进——从早期的 8-fact v1 基线到当前与运行时对齐的 13-fact v2——并在抽象层次混合和假阳性风险方面积累了经验教训。我们实现了一个完整的适配流水线：DCS 中的截图捕获、由大型 VLM 进行预标注、在 Label Studio 中进行人工审查、双语（英文和中文）监督微调导出、使用 Unsloth、PEFT 和 TRL 进行 LoRA 微调，以及在经核实训练数据零重叠的独立留出集上进行自动化基础模型与 LoRA 对比基准测试。此流水线中的每个工件——原始截图、预标注、人工审查标注、SFT 数据集、训练摘要、基准测试指标以及逐事实混淆图表——均存储在代码仓库中，可供独立检查和复现。此贡献在第\ref{ch:system-design}和\ref{ch:methodology}章中描述。
% 证据：C-VLM-01, C-VLM-02, C-VLM-03, C-VLM-03a, C-VLM-07, C-DATA-01, C-TRAIN-01

\item \textbf{基准测试结果表明小数据 LoRA 适配能显著提升座舱 VLM 在独立留出集上的事实提取准确率。}
在 13-fact v2 本体下，基于 928 条双语训练数据（Run-003 + Run-005$\times$2）训练的 Qwen3.5-9B LoRA 适配器在 Run-002 newfacts 留出集（50 样本）上实现事实准确率 0.9908（基础模型：0.8677），在 Run-004 random 留出集（100 样本）上实现 0.9931（基础模型：0.8038）。样本完全匹配率从 Run-002 的 0.10 升至 0.88，从 Run-004 的 0.03 升至 0.92。关键假阳性计数——即对下游具有直接风险、可能导致程序过早推进的错误——在 Run-002 上减少 64\%，在 Run-004 上减少 89\%。逐事实错误分析将完成类事实假阳性确定为主要残余失效模式。这些结果见第\ref{ch:evaluation}章，第\ref{sec:eval-tech}节。
% 证据：C-VLM-04, C-VLM-05

\item \textbf{骨干对比表明训练数据方案可跨模型架构泛化。}
相同的训练数据方案（Run-003 + Run-005$\times$2，928 条，双语）被应用于第二个 VLM 骨干 Gemma-4-31B。Gemma LoRA 适配器在 Run-002 newfacts 留出集上实现事实准确率 0.9492（基础模型：0.8169），在 Run-004 random 留出集上实现 0.9715（基础模型：0.8146），并在 Run-002 上实现零关键假阳性。虽然 Qwen 适配器在绝对准确率和召回率方面表现更优，但 Gemma 结果证实了改进是由训练数据方案——而非特定骨干的偶然产物——所驱动的，并揭示了骨干层面的精度-召回权衡差异，这对安全关键型部署具有重要参考价值。这些结果见第\ref{ch:evaluation}章，第\ref{sec:eval-gemma}节。
% 证据：C-VLM-06, C-VLM-03b

\item \textbf{针对未来 VR 人因研究的计划性评估框架。}
我们在设计层面规定了受控受试者间评估方案，比较视频加笔记、无状态基础纯文本 LLM 和基于状态证据教学三种条件，计划的评估指标涵盖即时表现（完成率、错误率、任务时间）、工作负荷（NASA-TLX）、学习效果（前/后测验）、保持（2--4 周延迟后测）和迁移（变体程序）。我们确定了从当前回放基础设施过渡到人因实验数据流水线所需的运行时数据采集升级，并编目了设计阶段已知的有效性威胁。此计划性评估在第\ref{ch:evaluation}章第\ref{sec:eval-planned}节中描述，代表研究计划的下一阶段。
% 证据：C-HUMAN-01, C-HUMAN-02, C-HUMAN-03, C-DATA-COLLECT-01（均为 TODO，被视为设计层面）
\end{enumerate}

贡献 1--4 由当前代码仓库中已有的已验证声明和基准测试工件支撑。贡献 5 是方法论和设计层面的贡献，记录了评估路径；它不报告人因数据，其所述指标和比较仍处于计划而非完成状态。

% ===========================================================================
\section{论文结构}\label{sec:intro-structure}

论文的剩余部分组织如下。

\textbf{第\ref{ch:background}}章提供理解本系统所需的领域和技术背景。它详细描述了 F/A-18C 冷启动程序，介绍了视觉事实本体及其从 8-fact v1 到 13-fact v2 的演进，并综述了智能教学系统、基于仿真器的培训、基于规则和状态感知的教学、面向结构化输出的 VLM 适配以及 AR/VR 环境中可解释 AI 辅助等相关工作。

\textbf{第\ref{ch:system-design}}章介绍 SimTutor 的设计。它从问题陈述中推导设计目标，描述端口-适配器架构，追踪完整帮助循环中的运行时数据流，并详细阐述遥测处理流水线、程序和门控规则引擎、仿真器适配层、事件日志与回放基础设施以及可选 VLM 适配组件的设计。该章最后与原论文研究计划进行系统性比较，记录设计在实现过程中的演进。

\textbf{第\ref{ch:methodology}}章描述 VLM 适配流水线的方法论框架。涵盖数据集构建（Run-001 至 Run-005，本体版本，标注分布，经核实零重叠的留出集构建），截图-预标注-审查-导出-训练-基准测试流水线，Qwen3.5-9B 和 Gemma-4-31B 的 LoRA 微调配置，以及基准测试协议，包括指标定义和污染控制。

\textbf{第\ref{ch:implementation}}章记录技术栈、核心领域逻辑模块、遥测与 DCS 集成、结构化帮助生成、VLM 视觉事实提取、回放与日志基础设施以及 CLI 和运行时入口点的实现细节。

\textbf{第\ref{ch:evaluation}}章分为两部分。A 部分报告当前技术结果：数据集和训练摘要、8-fact v1 基线、主要 13-fact v2 Qwen 基准测试结果、Gemma 骨干对比、逐事实错误分析，以及运行时基础设施就绪度评估。B 部分描述计划中的人因评估：原始评估愿景、计划指标、所需基础设施升级、拟议实验设计，以及已知有效性威胁编目。B 部分中的所有陈述均代表未来工作。

\textbf{第\ref{ch:discussion}}章对结果及其局限性进行解读。它考察了项目为何从 VR 优先转向遥测/VLM 重度实现，阐明当前基准测试能够和不能够证明的内容，分析本体从 8-fact 到 13-fact 的演进，论证人因证据为何仍然不可或缺，并识别将基准测试成功解读为学习成功的风险。

\textbf{第\ref{ch:conclusion}}章总结已完成的工作，重申贡献，并概述未来工作方向，包括 VR 运行时集成、人因评估、本体向其他程序的扩展、多帧时序推理以及在线适配。
